\documentclass[12pt]{article}

\usepackage{sbc-template}
\usepackage{graphicx}
\usepackage{url}
\usepackage{booktabs}
\usepackage{amsmath}
\usepackage{xcolor}
\usepackage{microtype}
\usepackage[hidelinks]{hyperref}
\usepackage{cite}
\usepackage{subcaption}

\newcommand{\ph}[1]{\colorbox{yellow!50}{\textbf{[#1]}}}

\sloppy

\title{Benchmarking Motif Discovery Models for Human ChIP-seq Data:\\
OOPS, ZOOPS, and STREME Applied to CTCF}

\author{}
\address{}

\begin{document}

\maketitle

\begin{abstract}
Motif discovery is a critical step in ChIP-seq data analysis, yet the statistical
model governing motif occurrence per sequence is rarely evaluated empirically.
This paper benchmarks three approaches --- MEME with the One Occurrence Per Sequence
(OOPS) model, MEME with the Zero Or One Occurrence Per Sequence (ZOOPS) model,
and STREME --- applied to CTCF ChIP-seq data from the \textit{H.\,sapiens} Calu-3
cell line (ENCODE experiment ENCSR921ERP).
The pipeline employs Bowtie2~2.5.5, SAMtools~1.23.1, Picard~3.4.0,
MACS3~3.0.4, and BEDTools~2.31.1.
From 36,355,553 raw reads, 35,775,372 passed quality filtering (98.4\%),
yielding 43,645 enrichment peaks at $q \leq 0.01$.
An order-2 Markov background model estimated from peak sequences was incorporated
to correct GC-content biases in regulatory regions.
ZOOPS improved the TOMTOM E-value by more than 10 orders of magnitude relative
to OOPS ($6.20 \times 10^{-20}$ vs.\ $3.37 \times 10^{-9}$), while the Markov
background alone had negligible effect on MEME, identifying the occurrence model
as the dominant source of PWM degradation.
Peak stratification confirmed that STREME recovers the CTCF motif across all
confidence strata, with statistical power increasing monotonically with sample
size.
All scripts are publicly available at \ph{anonymous repo URL}.

\textbf{Keywords:} ChIP-seq, CTCF, motif discovery, MEME, STREME, benchmarking,
transcription factor binding sites, PWM, GRCh38.
\end{abstract}

\section{Introduction}

Chromatin immunoprecipitation followed by sequencing (ChIP-seq) is a cornerstone
technique for mapping protein--DNA interactions genome-wide~\cite{park2009}.
After reads are aligned and enriched regions (peaks) are identified, motif
discovery determines the consensus DNA sequence recognized by the transcription
factor (TF) of interest~\cite{bailey2015}. This step relies on probabilistic
models that make explicit assumptions about how many times a motif occurs per
input sequence --- assumptions that directly affect the quality of the resulting
Position Weight Matrix (PWM).

Despite extensive ChIP-seq best-practice literature, the choice of occurrence
model in MEME is rarely evaluated empirically. The One Occurrence Per Sequence
(OOPS) model forces the algorithm to fit a motif to every submitted sequence,
which is biologically inconsistent with ChIP-seq data: low-confidence peaks,
experimental noise, and indirect binding produce sequences that do not contain
the target motif~\cite{bailey1995}. The Zero Or One Occurrence Per Sequence
(ZOOPS) model relaxes this constraint, allowing sequences without a motif to
contribute zero occurrences. STREME replaces the expectation-maximization (EM)
search entirely with suffix-tree enumeration and Fisher's exact test, scaling
naturally to large datasets without subsampling~\cite{bailey2021}.

This paper makes four contributions: (i)~a systematic comparison of OOPS,
ZOOPS, and STREME on the same ChIP-seq dataset, quantifying the effect of model
choice on PWM specificity and TOMTOM E-value; (ii)~an order-2 Markov background
model to correct nucleotide composition biases intrinsic to regulatory regions;
(iii)~a peak stratification analysis providing an empirical quality curve for
motif recovery as a function of peak confidence; and (iv)~a reproducible,
modular Bash pipeline with version-pinned tools, publicly available on a code
repository.

\section{Related Work}

Peak calling in ChIP-seq is predominantly performed with MACS~\cite{zhang2008};
successive releases (MACS2, MACS3) introduced improved FDR control via
Benjamini--Hochberg correction~\cite{feng2012}. HOMER~\cite{heinz2010} provides
an alternative caller with integrated motif discovery capabilities. Quality
control guidelines for ChIP-seq experiments, including the use of input controls
and ENCODE standards, are established in~\cite{landt2012}.

CTCF is one of the most-studied TFs for its role in loop extrusion and
topologically associating domain (TAD) organization~\cite{davidson2023,sun2022}.
Its 11-zinc-finger domain produces a highly asymmetric 19--20~bp motif whose
orientation determines whether two CTCF sites form convergent loop anchors.

Motif discovery algorithms range from EM-based MEME~\cite{bailey1995} to
the k-mer suffix-tree approach of STREME~\cite{bailey2021}. Downstream tools
include TOMTOM for PWM comparison~\cite{gupta2007} and FIMO for genome-wide
scanning~\cite{fimo2011}. Reference PWMs are provided by JASPAR~\cite{jaspar2024}
and HOCOMOCO~\cite{hocomoco2013}. Alternative motif discovery approaches
include BIOMAPP-CHIP and its recent BIOMAP extension, which employ k-mer mixture
models for ChIP-seq motif identification~\cite{caldonazzo2018,biomapp2024}.

Several integrated pipelines exist --- notably nf-core/chipseq (Nextflow-based)
and the ENCODE uniform processing pipeline --- but none provides an explicit,
quantitative comparison of occurrence models as the primary analytical
contribution. The present work fills this gap for CTCF in Calu-3, a lung
adenocarcinoma cell line underrepresented relative to reference lines K562
and GM12878.

\section{Case Study}

We analyzed ENCODE experiment \textbf{ENCSR921ERP}
(\url{https://www.encodeproject.org/experiments/ENCSR921ERP/}), profiling CTCF
in \textit{H.\,sapiens} Calu-3 cells. Data were generated on the Illumina
NextSeq~500 as 76-nt single-end reads (SE76). Sixteen ChIP FASTQ files
(libraries ENCLB483NAC and ENCLB214ERW) and twenty control (input) FASTQ files
(libraries ENCLB030UCW and ENCLB489EVC) were obtained from the ENCODE portal.

Calu-3 provides a biologically relevant context for studying CTCF in lung
cancer. The high-coverage dataset (36.3~M ChIP reads; 92.7~M control reads)
offers sufficient depth to evaluate motif recovery across different discovery
strategies and peak confidence strata.

\section{Material and Methods}

\subsection{Pipeline Architecture}

The pipeline comprises eleven numbered Bash scripts (\texttt{00\_config.sh}
through \texttt{10\_tomtom.sh}), with all parameters centralized in a single
configuration file (\texttt{00\_config.sh}). Tool versions are pinned via a
\texttt{micromamba} environment (\texttt{environment.yml}). All scripts and
the environment specification are available at \ph{anonymous repo URL}.

\subsection{Quality Control and Adapter Trimming}

Raw reads were assessed with \textbf{FastQC~0.12.1}. Adapters and low-quality
bases were removed by \textbf{Trim~Galore~2.2.0} with \texttt{--quality~20}
(Phred threshold) and \texttt{--length~36} (minimum post-trimming read length).

\subsection{Genome Alignment}

The GRCh38 primary assembly (Ensembl release~112) was indexed with
\textbf{Bowtie2~2.5.5}. Bowtie2 was chosen over Bowtie~1 for two reasons:
(1)~the 76-nt reads exceed Bowtie~1's optimal range ($\leq$50~nt); (2)~Bowtie2
supports gapped alignments, recovering reads discarded under the Bowtie~1 model.
Alignment parameters: \texttt{--very-sensitive -N~1 -L~20}; only primary
alignments with MAPQ $\geq 30$ were retained via \textbf{SAMtools~1.23.1}.

\subsection{Duplicate Marking}

PCR duplicates were flagged (but not removed) with \textbf{Picard~3.4.0}
\texttt{MarkDuplicates} (\texttt{REMOVE\_DUPLICATES=false},
\texttt{VALIDATION\_STRINGENCY=LENIENT}), preserving reads for MACS3's
internal duplicate estimation.

\subsection{Peak Calling}

Enrichment peaks were identified with \textbf{MACS3~3.0.4} \texttt{callpeak}
(\texttt{-g hs}, \texttt{-q 0.01}, \texttt{--call-summits}). The
$-10\log_{10}(q)$ score in the \texttt{narrowPeak} output was used for
downstream peak stratification.

\subsection{Sequence Extraction}

For each summit, a 200-bp window ($\pm$100~bp) was extracted from GRCh38 with
\textbf{BEDTools~2.31.1} (\texttt{slop} followed by \texttt{getfasta}).
Sequences with $>$10\% ambiguous bases (N) were discarded.

\subsection{Nucleotide Background Model}

Regulatory regions bound by CTCF --- CpG islands and promoters --- exhibit GC
content of 50--60\%. Under a uniform 25\% background, EM-based algorithms
over-represent GC-rich patterns unrelated to the true binding motif. An order-2
Markov model was estimated from the peak sequences using \texttt{fasta-get-markov}
(MEME Suite~5.5.9) and supplied via the \texttt{-bfile} parameter to all MEME
and STREME runs.

\subsection{Motif Discovery: Four Configurations}

All runs searched for one motif (\texttt{-nmotifs~1}) of width 19~nt on both
strands (\texttt{-revcomp}). Because MEME's documentation recommends limiting
input to $\sim$1,000 sequences, MEME runs used a random subset
(\texttt{sequences\_meme\_subset.fa}); STREME used the full peak set.

\begin{enumerate}
    \item \textbf{MEME OOPS} --- \texttt{-mod oops}, uniform background.
          Replicates the original pipeline as the explicit \emph{baseline}.
    \item \textbf{MEME OOPS + Markov} --- \texttt{-mod oops}, order-2 Markov.
          Isolates the effect of background correction independently.
    \item \textbf{MEME ZOOPS} --- \texttt{-mod zoops}, order-2 Markov.
          Allows noise-dominated sequences to contribute zero occurrences.
    \item \textbf{STREME} --- Suffix-tree enumeration with Fisher's exact test
          and order-2 Markov background, applied to the full peak set
          without subsampling.
\end{enumerate}

\subsection{Peak Stratification}

Peaks were sorted by MACS3 score and partitioned into four subsets: top 10\%
($n = 4{,}364$), top 25\% ($n = 10{,}911$), top 50\% ($n = 21{,}822$),
and all peaks ($n = 43{,}645$). STREME with Markov background was applied to
each subset independently.

\subsection{PWM Comparison}

All discovered motifs were compared against the JASPAR~2024 CORE
non-redundant database~\cite{jaspar2024} using \textbf{TOMTOM}~\cite{gupta2007}
with Pearson correlation distance (\texttt{-dist~pearson}), minimum overlap
of 5~nt, and significance threshold $p \leq 0.05$.

\section{Results and Discussion}

\subsection{Quality Control and Alignment}

After concatenation of the 16 ChIP FASTQ files, 36,355,553 raw reads were
obtained. FastQC detected high base quality (Phred $> 30$) with residual
TruSeq adapters. After Trim~Galore filtering, 35,775,372 reads were retained
(98.4\%). Bowtie2 achieved an overall alignment rate of 98.77\% for ChIP and
98.68\% for control. Picard flagged 7.31\% of ChIP reads and 5.44\% of control
reads as PCR duplicates. Full alignment statistics are presented in
Table~\ref{tab:alignment}.

\begin{table}[htbp]
\centering
\caption{Alignment statistics for ChIP and control samples (ENCODE ENCSR921ERP).}
\label{tab:alignment}
\small
\begin{tabular}{lrr}
\toprule
\textbf{Metric} & \textbf{ChIP} & \textbf{Control} \\
\midrule
Raw reads                  & 36,355,553  & 92,726,595 \\
Reads after trimming       & 35,775,372  & 92,726,595 \\
Overall alignment rate     & 98.77\%     & 98.68\%    \\
Reads with MAPQ $\geq 30$  & 35,295,000  & 80,007,586 \\
PCR duplicates (\%)        & 7.31\%      & 5.44\%     \\
\bottomrule
\end{tabular}
\end{table}

\subsection{Peak Calling}

MACS3 identified \textbf{43,645 enrichment peaks} at $q \leq 0.01$ using
only the FDR threshold, consistent with MACS3's mutually exclusive handling
of $p$-value and $q$-value parameters. The estimated average fragment size
was 175~bp, consistent with the expected 200--600~bp fragmentation range
for CTCF ChIP-seq.

\subsection{Comparative Motif Discovery}

Table~\ref{tab:comparison} reports TOMTOM results for all four configurations.
CTCF (JASPAR MA0139.2) was the top match in every configuration, confirming
biological validity across all tested approaches. However, the statistical
significance varies by several orders of magnitude depending on the model chosen.

\begin{table}[htbp]
\centering
\caption{Motif discovery benchmark. TOMTOM searched JASPAR~2024 CORE with
Pearson distance and $p \leq 0.05$; top match in all cases: CTCF (MA0139.2).
MEME runs used a 1,000-sequence random subset; STREME used all 43,645 peaks.}
\label{tab:comparison}
\small
\begin{tabular}{llrrr}
\toprule
\textbf{Approach} & \textbf{Background}
  & \textbf{$p$-value} & \textbf{E-value} & \textbf{$q$-value} \\
\midrule
MEME OOPS (baseline) & Uniform
  & $1.43{\times}10^{-12}$ & $3.37{\times}10^{-9}$ & $6.73{\times}10^{-9}$ \\
MEME OOPS + Markov   & Markov order-2
  & $1.42{\times}10^{-12}$ & $3.34{\times}10^{-9}$ & $6.68{\times}10^{-9}$ \\
MEME ZOOPS           & Markov order-2
  & $2.64{\times}10^{-23}$ & $6.20{\times}10^{-20}$ & $1.24{\times}10^{-19}$ \\
STREME (all peaks)   & Markov order-2
  & $4.46{\times}10^{-13}$ & $1.05{\times}10^{-9}$  & $2.09{\times}10^{-9}$  \\
\midrule
STREME (top 10\%)    & Markov order-2
  & $2.64{\times}10^{-10}$ & $6.20{\times}10^{-7}$  & $1.24{\times}10^{-6}$  \\
STREME (top 25\%)    & Markov order-2
  & $9.79{\times}10^{-12}$ & $2.30{\times}10^{-8}$  & $4.59{\times}10^{-8}$  \\
STREME (top 50\%)    & Markov order-2
  & $6.72{\times}10^{-13}$ & $1.58{\times}10^{-9}$  & $3.15{\times}10^{-9}$  \\
\bottomrule
\end{tabular}
\end{table}

The results reveal two independent effects. First, switching from OOPS to ZOOPS
improved the TOMTOM E-value by more than \textbf{10 orders of magnitude}
($3.37 \times 10^{-9}$ to $6.20 \times 10^{-20}$) under otherwise identical
conditions. This confirms that the OOPS assumption --- every sequence contains
exactly one binding site --- introduces substantial noise into the PWM: with
43,645 peaks at $q \leq 0.01$, a non-negligible fraction of lower-confidence
peaks represents noise or indirect binding. Forcing MEME to fit a motif to
each of these sequences dilutes the position weight matrix.

Second, the order-2 Markov background had a negligible effect on MEME OOPS
alone ($3.37 \times 10^{-9}$ vs.\ $3.34 \times 10^{-9}$), demonstrating that
the occurrence model is the dominant source of PWM degradation --- not background
composition bias. This is a practically important finding: correcting the
background without addressing the occurrence assumption provides essentially
no benefit.

STREME (all peaks, E-value $1.05 \times 10^{-9}$) is competitive with the
baseline OOPS but does not reach ZOOPS-level specificity in TOMTOM, despite
using all 43,645 sequences. This is expected: STREME's Fisher's exact test
optimizes for motif enrichment relative to a shuffled background rather than
for alignment fidelity to a reference PWM. The sequence logos generated by each
configuration are presented in Fig.~\ref{fig:logos}.

\begin{figure}[htbp]
    \centering
    \begin{subfigure}[b]{0.48\textwidth}
        \includegraphics[width=\textwidth]{figures/logo_meme_oops.eps}
        \caption{MEME OOPS --- uniform background (baseline)}
    \end{subfigure}
    \hfill
    \begin{subfigure}[b]{0.48\textwidth}
        \includegraphics[width=\textwidth]{figures/logo_meme_oops_markov.eps}
        \caption{MEME OOPS + Markov background}
    \end{subfigure}

    \vspace{0.5em}

    \begin{subfigure}[b]{0.48\textwidth}
        \includegraphics[width=\textwidth]{figures/logo_meme_zoops.eps}
        \caption{MEME ZOOPS + Markov background}
    \end{subfigure}
    \hfill
    \begin{subfigure}[b]{0.48\textwidth}
        \ph{figures/logo\_streme.png\\
        (exportar do streme.html:\\
        results/motifs/streme/streme.html)}
        \caption{STREME + Markov background (all peaks)}
    \end{subfigure}
    \caption{Sequence logos of the CTCF motif identified by each approach.
    Higher information content (bits) at conserved positions reflects greater
    PWM specificity. ZOOPS (c) shows markedly higher conservation at the
    central zinc-finger contact positions compared to OOPS (a).}
    \label{fig:logos}
\end{figure}

\subsection{Peak Stratification}

Table~\ref{tab:stratification} reports STREME motif E-values across four peak
confidence subsets, ordered by MACS3 score quantile.

\begin{table}[htbp]
\centering
\caption{STREME motif recovery as a function of peak confidence (Markov
background; top TOMTOM match in all cases: CTCF, MA0139.2).
E-values correspond to STREME's internal enrichment statistic.}
\label{tab:stratification}
\small
\begin{tabular}{lrrl}
\toprule
\textbf{Subset} & \textbf{Peaks ($n$)} & \textbf{STREME E-value} & \textbf{Top TOMTOM match} \\
\midrule
Top 10\%  & 4,364   & $5.1 \times 10^{-163}$ & CTCF (MA0139.2) \\
Top 25\%  & 10,911  & $3.8 \times 10^{-380}$ & CTCF (MA0139.2) \\
Top 50\%  & 21,822  & $5.9 \times 10^{-697}$ & CTCF (MA0139.2) \\
All peaks & 43,645  & $1.8 \times 10^{-933}$ & CTCF (MA0139.2) \\
\bottomrule
\end{tabular}
\end{table}

The STREME E-value improves monotonically as more peaks are included in the
analysis ($5.1 \times 10^{-163}$ for top~10\% to $1.8 \times 10^{-933}$ for
all peaks). This behavior reflects STREME's statistical framework: the Fisher's
exact test gains power as sample size increases, and the CTCF signal is
sufficiently strong across all confidence strata that lower-confidence peaks
contribute to statistical sensitivity rather than degrading the motif.

This finding has a direct implication for the OOPS failure mode. MEME OOPS
also receives more sequences as thresholds are relaxed, but cannot exclude
sequences without the motif. As a result, MEME OOPS degrades under large,
heterogeneous peak sets precisely where STREME improves. The stratification
analysis thus provides empirical, quantitative evidence that the OOPS assumption
is the critical limitation in large ChIP-seq experiments.

\subsection{Upgrade from Bowtie~1 to Bowtie2}

The original pipeline used Bowtie~1, which is optimized for reads of $\leq$50~nt
and does not support gapped alignments. The SE76 reads in this experiment exceed
that limit. Bowtie2~2.5.5 achieved a 98.77\% alignment rate for ChIP and 98.68\%
for control, confirming suitability for these read lengths. A direct head-to-head
comparison with Bowtie~1 statistics is proposed as future work.

\section{Conclusions}

We presented a systematic benchmarking of motif discovery models for ChIP-seq
data, demonstrating that the OOPS occurrence assumption --- prevalent in published
pipelines --- is inappropriate for large, heterogeneous peak sets. Switching from
OOPS to ZOOPS improved TOMTOM E-value by over 10 orders of magnitude, while the
Markov background correction alone had negligible impact on MEME, cleanly
separating the contributions of occurrence model and background model to PWM
quality. Peak stratification confirmed that STREME scales gracefully with sample
size, recovering the CTCF motif across all confidence strata with increasing
statistical power.

The comparative framework introduced here is broadly applicable: any ChIP-seq
study using MEME OOPS can benefit from switching to ZOOPS or STREME, particularly
for large peak sets. All analysis scripts are version-controlled and publicly
available, enabling full reproducibility.

Future work will: (i)~formalize the pipeline in Nextflow for cloud-portable
execution and nf-core compatibility; (ii)~apply the benchmarking framework to
additional TFs with distinct binding signatures to validate generalizability;
(iii)~extend the analysis with CENTRIMO for motif centrality assessment and AME
for co-factor enrichment in the Calu-3 context; and (iv)~profile CTCF site
orientation with FIMO to characterize TAD boundary topology in lung adenocarcinoma.

\bibliographystyle{sbc}
\bibliography{bibliography}

\end{document}
